\subsection{Hardware-aware activation-function choice}
\label{sec:activation-choice}

Models A and B use the rectified linear unit (ReLU),
$f(x)=\max(0,x)$, after every convolutional stage and hidden fully connected
layer.  This follows the course treatment in \emph{DNNs.pdf}: p.~97 describes
ReLU as computationally efficient ``thresholding'' and notes that it does not
saturate for positive inputs, while p.~98 recommends it for most cases.  The
same slide also identifies its limitations---negative units have zero gradient
and can become ``dying ReLUs'', and the output is not zero-centred.  These are
training risks, but neither changes the inference operation: compare with zero
and retain either zero or the input value.

The hardware consequence is substantial because activation cost is paid for
every feature value, not once per layer.  From the implemented output shapes,
each custom model evaluates
\begin{equation}
  32(62^2)+64(29^2)+128(12^2)+128+64=195{,}456
\end{equation}
hidden activation values per image.  Arm's CMSIS-NN Q7 reference ReLU implements
each scalar as a negative comparison followed by assignment of zero; on
DSP-capable M-profile cores its vector path masks four signed 8-bit values per
loop iteration.\footnote{Arm Ltd., \emph{CMSIS-NN
\texttt{arm\_relu\_q7.c}}, lines 47--94,
\url{https://github.com/ARM-software/CMSIS-NN/blob/main/Source/ActivationFunctions/arm_relu_q7.c}.}
Thus ``one comparison'' describes the logical arithmetic, although real cycle
counts also include loads, stores and loop control and depend on the core and
compiler.  ReLU needs no coefficients, transcendental unit or lookup table.

Sigmoid and tanh are less attractive here for both learning and deployment.
Their definitions,
\begin{equation}
  \sigma(x)=\frac{1}{1+e^{-x}}, \qquad
  \tanh(x)=\frac{e^x-e^{-x}}{e^x+e^{-x}},
\end{equation}
require exponential evaluation or an approximation.  \emph{DNNs.pdf}
pp.~95--96 also shows that both saturate, producing near-zero gradients at the
extremes.  A representative embedded alternative illustrates the memory--time
trade-off: CMSIS-NN implements 16-bit sigmoid and tanh by two reads from a
256-entry \texttt{uint16\_t} sigmoid table plus integer interpolation for each
value.  The table alone occupies $256\times2=512$ bytes of read-only memory,
before code and temporary values.\footnote{Arm Ltd., \emph{CMSIS-NN
\texttt{arm\_nn\_activation\_s16.c}}, lines 47--108, and
\texttt{arm\_nn\_tables.h}, line 37,
\url{https://github.com/ARM-software/CMSIS-NN/blob/main/Source/ActivationFunctions/arm_nn_activation_s16.c},
\url{https://github.com/ARM-software/CMSIS-NN/blob/main/Include/arm_nn_tables.h}.}
Using either function at all hidden sites would therefore repeat table access
and interpolation 195,456 times per image instead of applying a threshold.

Floating-point support cannot be assumed on a small microcontroller.  Arm's
Cortex-M DSP study reports that Cortex-M3 executes floating-point work by
integer emulation, while FPU-equipped Cortex-M4/M7 devices accelerate it; in
the paper's 32-point floating-point real-FFT benchmark, the measured counts are
16,183 cycles on Cortex-M3 and 962 on Cortex-M7, a $16.8\times$ difference.
The same paper summarises the broader measured gap as approximately
$20\times$.\footnote{Arm Ltd., \emph{DSP capabilities of Cortex-M4 and
Cortex-M7}, 2016, pp.~3 and 12,
\url{https://developer.arm.com/community/arm-community-blogs/b/architectures-and-processors-blog/posts/white-paper-dsp-capabilities-of-cortex-m4-and-cortex-m7}.}
This is not an \texttt{exp} benchmark, so it is not used as a universal
activation latency; it demonstrates why an FPU-specific cycle claim would be
misleading and why avoiding transcendental floating-point work is the safer
architecture choice.

GELU was also rejected.  Its exact form $x\Phi(x)$ uses the Gaussian cumulative
distribution (equivalently an error function), while its common approximation
uses tanh; either is more involved than thresholding.\footnote{D. Hendrycks and
K. Gimpel, \emph{Gaussian Error Linear Units (GELUs)}, 2016,
\url{https://arxiv.org/abs/1606.08415}.}
Leaky ReLU and PReLU are computationally efficient alternatives discussed on
p.~99 of \emph{DNNs.pdf} and avoid dead neurons, but they require a multiply on
the negative branch; PReLU additionally stores and learns the slope $\alpha$.
Without evidence that dying ReLUs impair these models, those extra operations
and parameters are not justified.

\subsubsection{ReLU6 and low-precision deployment}

MobileNetV2 uses
$\operatorname{ReLU6}(x)=\min(\max(0,x),6)$ because the bounded activation is
robust under low-precision computation.\footnote{M. Sandler et al.,
\emph{MobileNetV2: Inverted Residuals and Linear Bottlenecks}, Section~4,
\url{https://arxiv.org/abs/1801.04381}.}
The upper bound prevents rare large activations from expanding the represented
range, so a fixed number of integer levels can cover the useful interval more
densely.  This is directly relevant to int8 deployment and explains why Arm
CMSIS-NN provides a dedicated signed-8-bit ReLU6 operator.  However, this
assignment makes the edge argument analytically and explicitly excludes
quantisation.  The custom models are trained and evaluated in floating point,
so imposing an upper saturation at six would discard positive magnitude
without delivering a required int8 benefit.  Ordinary ReLU therefore matches
both the lecture recommendation and the present deployment scope; ReLU6 is a
future quantised-deployment option rather than the chosen activation.

\subsubsection{The softmax exception}

The ten-class output is the one place where exponentials are justified when
normalised probabilities are required.  As shown in \emph{Classification
2.pdf}, pp.~4--5,
\begin{equation}
  p(y=C_k\mid x)=\frac{\exp(z_k)}{\sum_{j=1}^{10}\exp(z_j)}.
\end{equation}
This softmax is evaluated once over ten logits: ten exponentials per image,
not once over a feature map.  Replacing all hidden ReLUs by an exponential-based
activation would apply the expensive operation to 195,456 values, approximately
19,546 times as many values as the output softmax.

The PyTorch models deliberately return raw logits.  During training,
\texttt{CrossEntropyLoss} consumes those logits directly, avoiding a separate
softmax layer.  During inference, softmax is needed only when calibrated class
probabilities must be reported; the predicted class itself is the index of the
largest logit because softmax preserves ordering.  Consequently, the necessary
probability normalisation is confined to a negligible ten-value output stage,
while the high-volume hidden activations retain the comparison-only ReLU.
